\subsection{RQ7: How do SLMs compare with substantially larger language models across code generation, error detection, and tool-assisted settings?}
\label{sec:rq7}

This RQ synthesizes how the evaluated SLMs compare with two substantially larger language models across code generation, error detection, specialization, static-analysis-assisted detection, and iterative generation. All models are evaluated on the same EvalPlus benchmark and under the same execution-based reference used throughout the study whenever the corresponding experimental condition permits a direct comparison. This is not a causal scale experiment. The two larger models were purposively selected, and parameter count varies together with model family, training, architecture, backend, provider revision, and numerical precision. The results therefore characterize the evaluated models under a common experimental framework rather than isolate the effect of model scale.

\subsubsection{Code Generation}

The code-generation experiment provides the clearest evidence that the substantially larger models do not establish a separate upper performance range on these tasks. The three highest observed proportions of \emph{Correct} solutions are all obtained by SLMs under the same EvalPlus evaluation harness: Muse Glimmer reaches 80.8\%, GPT-OSS 79.0\%, and Gemma 4 31B 78.7\%. DeepSeek V4.1 Flash and MiMo V2.5 each reach 73.4\%, which is also the result obtained by Qwen 2.5 Coder 7B.

Thus, selected SLMs do not merely fall within the same observed range as the larger baselines. Several obtain higher correctness proportions under the evaluated generation protocol despite a substantially smaller nominal parameter count. The strongest SLMs in this comparison contain approximately 20B to 31B parameters, whereas the two hosted baselines contain 310B and 552B parameters. These results do not establish that smaller models are generally superior, but they show that an approximately order-of-magnitude increase in parameter count does not guarantee a higher Correct proportion on the evaluated function-level tasks.

The absence of a simple scale ordering also appears within a single model family. Qwen 2.5 Coder 7B reaches 73.4\% Correct, compared with 67.1\% for the 32B configuration. This observation is consistent with the cross-family comparison: increasing parameter count alone is insufficient to explain the differences observed under this generation protocol.

Larger scale also does not eliminate structural or output-contract failures. DeepSeek produces approximately 5.9\% Syntax outcomes, largely associated with responses that do not satisfy the expected function-level output structure. Most SLMs remain below approximately 3\%, although Ornith is an important exception with a substantially higher Syntax rate. The comparison therefore indicates that instruction following and compliance with the expected generation interface remain relevant even for substantially larger models.

\subsubsection{Hallucination Detection}

The hallucination-detection results show a similar absence of a distinct performance range for the larger models. Muse Glimmer reaches 73.1\% Accuracy and 53.8\% Macro-F1 among the general-purpose SLMs. DeepSeek reaches 75.5\% Accuracy and 51.0\% Macro-F1, while MiMo reaches 51.5\% Accuracy and 39.7\% Macro-F1. Muse therefore approaches DeepSeek in Accuracy while exceeding both larger baselines in Macro-F1.

The comparison becomes more pronounced after specialization. Best Acc reaches 77.8\% Accuracy, while Weighted MF1 reaches 54.2\% Macro-F1. Both values exceed the corresponding aggregate results of the two general-purpose larger models. These improvements are obtained without increasing the parameter count of the underlying Gemma model, indicating that task specialization can alter detection capability substantially without relying on additional model scale.

The class-specific results also show that the main limitation identified in RQ2 persists for the substantially larger models. \emph{Runtime Error} remains difficult across both groups. DeepSeek reaches 14.8\% Runtime Recall and MiMo reaches 21.5\%, while GPT-OSS reaches 32.9\% and several other SLM configurations also fall within or above the large-model range. Thus, increasing scale from tens to hundreds of billions of parameters does not remove the difficulty of recognizing Runtime failures under the evaluated protocol.

These results reinforce the distinction between generation and evaluation capability. Some SLMs are strong generators but weak detectors, while specialized SLMs improve detection substantially without increasing model size. The substantially larger baselines exhibit the same broader separation between these capabilities, indicating that the difficulty of recognizing generated-code failures is not resolved by parameter scale alone.

\subsubsection{Tool-Assisted Error Detection}

Static-analysis feedback likewise does not produce a clear advantage for the substantially larger models. The strongest effect remains the improvement in \emph{Syntax Error} detection. Most SLMs reach 100\% or near-100\% Syntax Recall after receiving compilation feedback, and both larger models reach 100\%. Once deterministic syntactic evidence is supplied explicitly, the evaluated SLMs are therefore generally able to exploit it as effectively as the substantially larger baselines.

The aggregate results further show that selected SLMs remain highly competitive after external evidence is introduced. Muse Glimmer reaches approximately 79\% Accuracy and 71\% Macro-F1, while Gemma 4 31B reaches 77\% Accuracy and 66\% Macro-F1. The specialized Weighted MF1 detector reaches approximately 76\% Accuracy and 68\% Macro-F1. DeepSeek reaches approximately 76\% Accuracy and 64\% Macro-F1, while MiMo reaches 47\% for both metrics. Thus, static feedback does not create a new performance separation in favor of the substantially larger models. Several SLM conditions equal or exceed the observed aggregate results of the hosted baselines.

The Runtime limitation also persists after tool integration. DeepSeek decreases from 14.8\% to approximately 10\% Runtime Recall, while MiMo decreases from 21.5\% to approximately 9\%. Similar reductions occur for many SLMs. The difficulty identified in RQ5 is therefore shared across the evaluated scale range: incomplete Pyright evidence is not reliably translated into an explicit \emph{Runtime Error} prediction even by models with hundreds of billions of parameters.

The indirect effect of static feedback on classes not evaluated by the tools is also present at larger scale. DeepSeek decreases Correct Recall while increasing Functional Recall, and MiMo presents an even clearer redistribution, with Correct Recall decreasing by approximately 10 percentage points while Functional Recall increases by approximately 10 points. The Correct-to-Functional shifts observed among several SLMs in RQ5 are therefore not exclusive to compact models. Both groups can interpret static diagnostics as broader evidence about program correctness, even when the tools provide no direct information about functional behavior.

Together, these results suggest that the effect of tool assistance depends more strongly on the type and coverage of the supplied evidence and on the model's learned decision behavior than on parameter count alone. Direct compilation evidence is successfully incorporated across model scales, whereas partial Runtime evidence remains difficult to interpret consistently.

\subsubsection{Tool-Assisted Iterative Code Generation}

The iterative generation experiment presents a somewhat different pattern. Among the eleven SLMs, six increase their proportion of \emph{Correct} solutions relative to RQ1, while five decrease it. In contrast, both substantially larger models improve. DeepSeek increases Correct by 7.3 percentage points and MiMo by 6.6 points. RQ6 provides the corresponding paired analysis for these two models.

This difference indicates a more consistently positive response to the iterative protocol for the two larger models evaluated in this study. However, the improvement does not establish a higher final correctness range. Gemma 4 31B reaches approximately 82\% Correct after the iterative protocol, compared with approximately 81\% for DeepSeek and 80\% for MiMo. GPT-OSS also reaches approximately 80\%. The highest final Correct proportion therefore remains associated with an SLM despite the more consistent direction of change observed for the two substantially larger models.

The error distributions also reveal similarities across scales. Runtime and Syntax failures frequently decrease during iterative generation for both SLMs and the larger models. Functional behavior remains more model-dependent, and several SLMs achieve correctness gains through reductions that extend beyond the error categories directly targeted by the tools. Qwen 2.5 Coder 32B, for example, obtains the largest SLM increase in Correct while substantially reducing Functional failures, despite the absence of direct functional evidence from compilation or Pyright.

The more consistent improvement observed for DeepSeek and MiMo therefore appears to concern the response to the complete iterative protocol rather than the maximum final performance reached. With only two substantially larger models, this pattern is insufficient to establish that model scale causes more reliable use of feedback. It instead motivates the narrower observation that both hosted baselines respond positively in this experiment, whereas the SLM responses are divided between improvement and regression.

Across the evaluated settings, increasing parameter count does not produce a uniform performance advantage. Under the same benchmark and execution-based evaluation framework, the three highest code-generation Correct proportions are obtained by SLMs, specialized SLMs reach or exceed the larger baselines in hallucination detection, and selected SLMs also equal or surpass their aggregate performance when static-analysis evidence is provided. Runtime detection remains difficult at both scales, and the Correct-to-Functional redistribution caused by tool feedback also appears in the substantially larger models. The iterative protocol is the main setting in which the two larger baselines exhibit a more consistently positive direction of change, yet their final correctness remains within the range reached by the strongest SLMs. The results therefore do not support a general performance ordering based on parameter count and instead show that specialization, error type, and the ability to interpret and act on external evidence are important dimensions of model reliability under the evaluated protocols.